% 逐表替换用：不要和已内嵌这些表格的正文重复载入。
\urlstyle{same}

% ===== tab:p1_acoustic_groups =====
\begin{table}[H]
    \centering\small
    \caption{25维声学特征的组成及含义}
    \label{tab:p1_acoustic_groups}
    \setlength{\tabcolsep}{6pt}
    \renewcommand{\arraystretch}{1.12}
    \ifdefined\pOneCompactTableBox\else\newsavebox{\pOneCompactTableBox}\fi
    \begin{lrbox}{\pOneCompactTableBox}
    \begin{tabular}{l>{\raggedright\arraybackslash}p{17em}>{\raggedright\arraybackslash}p{10em}c}
        \toprule
        类别 & 描述量 & 主要描述内容 & 维度\\
        \midrule
        响度 & 听觉加权与压缩后的Loudness & 感知强弱 & 1\\
        基频 & 相对于27.5 Hz的半音基频 & 音高变化 & 1\\
        频谱结构 & Alpha比、Hammarberg指数、两项谱斜率、谱通量 & 频谱形状及其变化 & 5\\
        倒谱 & 第1至第4阶MFCC & 谱包络的低阶表示 & 4\\
        音质 & Jitter、Shimmer的分贝表示、基于自相关估计的HNR、H1--H2、H1--A3 & 周期稳定性与谐波结构 & 5\\
        共振峰 & 前三阶共振峰的频率、带宽及相对于基频谐波的对数幅度 & 声道共振特性 & 9\\
        \midrule
        合计 & & & 25\\
        \bottomrule
    \end{tabular}
    \end{lrbox}
    \usebox{\pOneCompactTableBox}
\end{table}

% ===== tab:p1_acoustic_example =====
\begin{table}[H]
    \centering\small
    \caption{两组声学观测的数值与有效性对照}
    \label{tab:p1_acoustic_example}
    \setlength{\tabcolsep}{6pt}
    \renewcommand{\arraystretch}{1.12}
    \ifdefined\pOneCompactTableBox\else\newsavebox{\pOneCompactTableBox}\fi
    \begin{lrbox}{\pOneCompactTableBox}
    \begin{tabular}{cccccc}
        \toprule
        记录区间/s & 中心时刻/s & 响度 & 半音基频 & MFCC1 & 记录有效\\
        \midrule
        $[1.49,1.51]$ & 1.50 & 0.870 & 36.368 & 36.947 & 是\\
        $[1.99,2.01]$ & 2.00 & 0.363 & $0^{\ast}$ & $-4.455$ & 是\\
        \bottomrule
    \end{tabular}
    \end{lrbox}
    \usebox{\pOneCompactTableBox}
\par\smallskip
\parbox{\wd\pOneCompactTableBox}{\footnotesize 注：$0^{\ast}$表示未输出有声基频时保留的零标记，不作为实际音高解释。}
\end{table}

% ===== tab:p1_visual_groups =====
\begin{table}[H]
    \centering\small
    \caption{52维面部系数的分组及代表性含义}
    \label{tab:p1_visual_groups}
    \setlength{\tabcolsep}{6pt}
    \renewcommand{\arraystretch}{1.12}
    \ifdefined\pOneCompactTableBox\else\newsavebox{\pOneCompactTableBox}\fi
    \begin{lrbox}{\pOneCompactTableBox}
    \begin{tabular}{l>{\raggedright\arraybackslash}p{16em}>{\raggedright\arraybackslash}p{11em}c}
        \toprule
        分组 & 代表性输出 & 主要描述内容 & 维度\\
        \midrule
        额眉部 & 压眉、内侧抬眉、外侧抬眉 & 眉部位置与形态 & 5\\
        眼部 & 眨眼、睁大、眯眼及视线方向 & 眼睑状态与眼部运动 & 14\\
        面颊与鼻部 & 鼓腮、面颊收紧、鼻部牵动 & 面中部形态 & 5\\
        嘴部 & 嘴角上扬或下压、噘嘴、嘴唇收紧等 & 嘴唇及嘴角形态 & 23\\
        下颌 & 张嘴、下颌前伸及左右移动 & 下颌位置与开合 & 4\\
        其他输出项 & 工具字段\_neutral & 保留输出，不作为中性情感判据 & 1\\
        \midrule
        合计 & & & 52\\
        \bottomrule
    \end{tabular}
    \end{lrbox}
    \usebox{\pOneCompactTableBox}
\end{table}

% ===== tab:p1_result_overview =====
\begin{table}[H]
    \centering\small
    \caption{样本处理与三种模态的观测覆盖情况}
    \label{tab:p1_result_overview}
    \setlength{\tabcolsep}{6pt}
    \renewcommand{\arraystretch}{1.12}
    \ifdefined\pOneCompactTableBox\else\newsavebox{\pOneCompactTableBox}\fi
    \begin{lrbox}{\pOneCompactTableBox}
    \begin{tabular}{lrr}
        \toprule
        统计项目 & 数量 & 占比\\
        \midrule
        已处理样本 & 100 & 100.00\%\\
        与原文对应的输出词位置 & 1932 & 100.00\%\\
        文本模态可观测的词 & 1932 & 100.00\%\\
        声学模态可观测的词 & 1932 & 100.00\%\\
        视觉模态可观测的词 & 1595 & 82.56\%\\
        \bottomrule
    \end{tabular}
    \end{lrbox}
    \usebox{\pOneCompactTableBox}
\par\smallskip
\end{table}

% ===== tab:p1_manual_review =====
\begin{table}[H]
    \centering\small
    \caption{21条复核样本的人工审核结果}
    \label{tab:p1_manual_review}
    \setlength{\tabcolsep}{6pt}
    \renewcommand{\arraystretch}{1.12}
    \ifdefined\pOneCompactTableBox\else\newsavebox{\pOneCompactTableBox}\fi
    \begin{lrbox}{\pOneCompactTableBox}
    \begin{tabular}{lcc>{\raggedright\arraybackslash}p{18em}}
        \toprule
        审核发现 & 样本数 & 原文词数 & 主要表现\\
        \midrule
        音轨与文本不对应 & 11 & 120 & 实际语音内容与给定转写不对应\\
        无对应朗读 & 5 & 64 & 未听到对应朗读，其中1条为音乐\\
        非英语语音 & 3 & 93 & 实际语音并非英语\\
        分段内容不对应 & 1 & 18 & 前半段无朗读，后半段语音与转写不对应\\
        音频卡顿、局部对应 & 1 & 22 & 音频存在卡顿，部分词语与文本一致，但未形成完整对应\\
        \midrule
        合计 & 21 & 317 & \\
        \bottomrule
    \end{tabular}
    \end{lrbox}
    \usebox{\pOneCompactTableBox}
\end{table}

% ===== tab:p1_mapping_records =====
\begin{table}[H]
    \centering\small
    \caption{典型样本的部分词级映射记录}
    \label{tab:p1_mapping_records}
    \setlength{\tabcolsep}{6pt}
    \renewcommand{\arraystretch}{1.12}
    \ifdefined\pOneCompactTableBox\else\newsavebox{\pOneCompactTableBox}\fi
    \begin{lrbox}{\pOneCompactTableBox}
    \begin{tabular}{lcccc}
        \toprule
        原文词 & 子词索引 & 候选区间/s & 声学记录索引 & 视觉帧时刻/s\\
        \midrule
        started & 1 & $[1.2825,1.6425]$ & 127--164 & 1.2, 1.4, 1.6\\
        successful & 3 & $[1.8425,2.4025]$ & 183--240 & 1.8, 2.0, 2.2, 2.4\\
        Consulting & 6 & $[3.0225,3.5225]$ & 301--352 & 3.0, 3.2, 3.4, 3.6\\
        \bottomrule
    \end{tabular}
    \end{lrbox}
    \usebox{\pOneCompactTableBox}
\par\smallskip
\parbox{\wd\pOneCompactTableBox}{\footnotesize 注：索引从0开始；声学索引指具有正聚合权重的原生特征记录，不是波形采样点。视觉帧按其时间支持区间与候选词区间的交叠关系参与聚合，帧中心不必位于词区间内部。}
\end{table}

% ===== tab:p1_feature_specification =====
\begin{table}[H]
    \centering\small
    \caption{三种模态的词级特征规格}
    \label{tab:p1_feature_specification}
    \setlength{\tabcolsep}{6pt}
    \renewcommand{\arraystretch}{1.12}
    \ifdefined\pOneCompactTableBox\else\newsavebox{\pOneCompactTableBox}\fi
    \begin{lrbox}{\pOneCompactTableBox}
    \begin{tabular}{lccc}
        \toprule
        模态 & 特征维数 & 样本矩阵形状 & 对齐粒度\\
        \midrule
        文本 & 768 & $L\times768$ & 原文词\\
        声学 & 25 & $L\times25$ & 原文词\\
        视觉（面部） & 52 & $L\times52$ & 原文词\\
        \bottomrule
    \end{tabular}
    \end{lrbox}
    \usebox{\pOneCompactTableBox}
\par\smallskip
\parbox{\wd\pOneCompactTableBox}{\footnotesize 注：$L$为对应样本的原文词数。各样本按实际词数存储，填充数与截断词数均为0；视觉模态不可观测的词位置采用零向量占位，并附无效标记。}
\end{table}

% ===== tab:p1_all_samples =====
\begingroup
\small
\setlength{\tabcolsep}{6pt}
\renewcommand{\arraystretch}{1.12}
\setlength{\LTleft}{\fill}
\setlength{\LTright}{\fill}
\begin{longtable}{lrrrrrr}
\caption{附件1全部样本的特征输出与有效位置统计}\label{tab:p1_all_samples}\\
\toprule
样本编号 & \shortstack{片段\\时长/s} & \shortstack{音频\\时长/s} & \shortstack{原文\\词数} & \shortstack{视觉可观测\\词数} & \shortstack{通过对齐\\筛选词数} & \shortstack{其中三种模态\\都可观测\\词数}\\
\midrule
\endfirsthead
\multicolumn{7}{c}{表~\thetable\quad 特征输出与有效位置统计（续）}\\
\toprule
样本编号 & \shortstack{片段\\时长/s} & \shortstack{音频\\时长/s} & \shortstack{原文\\词数} & \shortstack{视觉可观测\\词数} & \shortstack{通过对齐\\筛选词数} & \shortstack{其中三种模态\\都可观测\\词数}\\
\midrule
\endhead
\midrule
\multicolumn{7}{r}{续下页}\\
\endfoot
\bottomrule
\endlastfoot
\multicolumn{7}{l}{\textbf{1. 全部词通过筛选且三种模态都可观测（32条）}}\\*
\nolinkurl{-3g5yACwYnA$_$3} & 14.389 & 14.328 & 29 & 29 & 29 & 29\\
\nolinkurl{-3nNcZdcdvU$_$5} & 7.868 & 7.789 & 18 & 18 & 18 & 18\\
\nolinkurl{-6rXp3zJ3kc$_$8} & 12.805 & 12.738 & 34 & 34 & 34 & 34\\
\nolinkurl{-I_e4mIh0yE$_$1} & 7.622 & 7.444 & 18 & 18 & 18 & 18\\
\nolinkurl{-I_e4mIh0yE$_$3} & 9.163 & 8.994 & 22 & 22 & 22 & 22\\
\nolinkurl{-MeTTeMJBNc$_$0} & 9.333 & 9.218 & 21 & 21 & 21 & 21\\
\nolinkurl{-MeTTeMJBNc$_$13} & 5.430 & 5.255 & 14 & 14 & 14 & 14\\
\nolinkurl{-RfYyzHpjk4$_$11} & 5.239 & 5.075 & 17 & 17 & 17 & 17\\
\nolinkurl{-RfYyzHpjk4$_$8} & 4.589 & 4.532 & 17 & 17 & 17 & 17\\
\nolinkurl{-UUCSKoHeMA$_$0} & 7.867 & 7.732 & 18 & 18 & 18 & 18\\
\nolinkurl{-UacrmKiTn4$_$10} & 7.362 & 7.242 & 18 & 18 & 18 & 18\\
\nolinkurl{-UacrmKiTn4$_$4} & 4.741 & 4.682 & 14 & 14 & 14 & 14\\
\nolinkurl{-aNfi7CP8vM$_$7} & 8.695 & 8.538 & 18 & 18 & 18 & 18\\
\nolinkurl{-aqamKhZ1Ec$_$0} & 11.000 & 10.890 & 14 & 14 & 14 & 14\\
\nolinkurl{-dxfTGcXJoc$_$1} & 16.161 & 16.018 & 37 & 37 & 37 & 37\\
\nolinkurl{-egA8-b7-3M$_$16} & 5.538 & 5.483 & 11 & 11 & 11 & 11\\
\nolinkurl{-egA8-b7-3M$_$17} & 7.271 & 7.095 & 16 & 16 & 16 & 16\\
\nolinkurl{-egA8-b7-3M$_$18} & 8.386 & 8.266 & 22 & 22 & 22 & 22\\
\nolinkurl{-egA8-b7-3M$_$20} & 6.645 & 6.474 & 20 & 20 & 20 & 20\\
\nolinkurl{-egA8-b7-3M$_$26} & 6.031 & 5.921 & 15 & 15 & 15 & 15\\
\nolinkurl{-egA8-b7-3M$_$6} & 6.265 & 6.183 & 12 & 12 & 12 & 12\\
\nolinkurl{-egA8-b7-3M$_$9} & 5.090 & 4.936 & 10 & 10 & 10 & 10\\
\nolinkurl{-mqbVkbCndg$_$0} & 6.500 & 6.386 & 14 & 14 & 14 & 14\\
\nolinkurl{-s9qJ7ATP7w$_$6} & 2.475 & 2.315 & 5 & 5 & 5 & 5\\
\nolinkurl{-s9qJ7ATP7w$_$8} & 3.295 & 3.184 & 12 & 12 & 12 & 12\\
\nolinkurl{-t217m2on-s$_$7} & 17.183 & 17.026 & 45 & 45 & 45 & 45\\
\nolinkurl{-tPCytz4rww$_$10} & 4.789 & 4.679 & 12 & 12 & 12 & 12\\
\nolinkurl{-tPCytz4rww$_$11} & 5.467 & 5.379 & 17 & 17 & 17 & 17\\
\nolinkurl{-tPCytz4rww$_$18} & 6.645 & 6.569 & 16 & 16 & 16 & 16\\
\nolinkurl{-uywlfIYOS8$_$4} & 6.045 & 5.881 & 18 & 18 & 18 & 18\\
\nolinkurl{-wny0OAz3g8$_$0} & 3.333 & 3.274 & 9 & 9 & 9 & 9\\
\nolinkurl{-wny0OAz3g8$_$2} & 6.687 & 6.523 & 23 & 23 & 23 & 23\\
\addlinespace[0.5em]
\multicolumn{7}{l}{\textbf{2. 部分词通过筛选且三种模态都可观测（46条）}}\\*
\nolinkurl{-3g5yACwYnA$_$13} & 5.514 & 5.399 & 15 & 15 & 14 & 14\\
\nolinkurl{-3g5yACwYnA$_$2} & 9.394 & 9.231 & 14 & 14 & 13 & 13\\
\nolinkurl{-3g5yACwYnA$_$9} & 8.817 & 8.696 & 23 & 23 & 21 & 21\\
\nolinkurl{-571d8cVauQ$_$0} & 5.167 & 4.992 & 12 & 12 & 10 & 10\\
\nolinkurl{-571d8cVauQ$_$5} & 15.273 & 15.115 & 43 & 43 & 42 & 42\\
\nolinkurl{-9y-fZ3swSY$_$0} & 6.900 & 6.757 & 22 & 22 & 19 & 19\\
\nolinkurl{-9y-fZ3swSY$_$4} & 2.821 & 2.665 & 9 & 9 & 6 & 6\\
\nolinkurl{-9y-fZ3swSY$_$8} & 4.989 & 4.917 & 12 & 3 & 11 & 2\\
\nolinkurl{-AUZQgSxyPQ$_$2} & 22.511 & 22.405 & 41 & 41 & 39 & 39\\
\nolinkurl{-HeZS2-Prhc$_$2} & 8.262 & 8.151 & 16 & 16 & 15 & 15\\
\nolinkurl{-HwX2H8Z4hY$_$2} & 3.982 & 3.916 & 7 & 1 & 4 & 1\\
\nolinkurl{-HwX2H8Z4hY$_$5} & 5.653 & 5.544 & 9 & 2 & 8 & 2\\
\nolinkurl{-HwX2H8Z4hY$_$6} & 3.358 & 3.256 & 7 & 5 & 6 & 5\\
\nolinkurl{-MeTTeMJBNc$_$7} & 10.517 & 10.411 & 31 & 31 & 29 & 29\\
\nolinkurl{-RfYyzHpjk4$_$2} & 5.516 & 5.392 & 25 & 24 & 21 & 21\\
\nolinkurl{-THoVjtIkeU$_$12} & 14.896 & 14.806 & 39 & 39 & 38 & 38\\
\nolinkurl{-THoVjtIkeU$_$2} & 4.292 & 4.231 & 10 & 10 & 9 & 9\\
\nolinkurl{-THoVjtIkeU$_$6} & 8.097 & 7.983 & 30 & 30 & 27 & 27\\
\nolinkurl{-UuX1xuaiiE$_$0} & 3.700 & 3.576 & 9 & 6 & 8 & 5\\
\nolinkurl{-UuX1xuaiiE$_$1} & 10.350 & 10.268 & 27 & 1 & 27 & 1\\
\nolinkurl{-UuX1xuaiiE$_$3} & 4.130 & 4.038 & 9 & 9 & 8 & 8\\
\nolinkurl{-UuX1xuaiiE$_$6} & 8.114 & 7.944 & 22 & 22 & 20 & 20\\
\nolinkurl{-a55Q6RWvTA$_$3} & 22.154 & 22.083 & 65 & 60 & 63 & 58\\
\nolinkurl{-dxfTGcXJoc$_$0} & 20.533 & 20.410 & 40 & 33 & 39 & 32\\
\nolinkurl{-dxfTGcXJoc$_$2} & 16.698 & 16.633 & 34 & 34 & 33 & 33\\
\nolinkurl{-dxfTGcXJoc$_$6} & 12.735 & 12.628 & 28 & 28 & 27 & 27\\
\nolinkurl{-egA8-b7-3M$_$1} & 10.266 & 10.156 & 20 & 20 & 18 & 18\\
\nolinkurl{-egA8-b7-3M$_$13} & 4.184 & 4.024 & 11 & 11 & 10 & 10\\
\nolinkurl{-lzEya4AM_4$_$5} & 7.675 & 7.570 & 19 & 19 & 18 & 18\\
\nolinkurl{-lzEya4AM_4$_$6} & 14.792 & 14.740 & 43 & 43 & 42 & 42\\
\nolinkurl{-qDkUB0GgYY$_$6} & 4.678 & 4.615 & 16 & 16 & 12 & 12\\
\nolinkurl{-s9qJ7ATP7w$_$0} & 6.867 & 6.711 & 24 & 24 & 23 & 23\\
\nolinkurl{-s9qJ7ATP7w$_$1} & 4.792 & 4.644 & 11 & 11 & 10 & 10\\
\nolinkurl{-s9qJ7ATP7w$_$4} & 4.427 & 4.247 & 9 & 6 & 8 & 5\\
\nolinkurl{-s9qJ7ATP7w$_$5} & 8.561 & 8.398 & 19 & 19 & 18 & 18\\
\nolinkurl{-s9qJ7ATP7w$_$7} & 6.966 & 6.839 & 28 & 14 & 24 & 12\\
\nolinkurl{-t217m2on-s$_$2} & 5.686 & 5.630 & 12 & 12 & 11 & 11\\
\nolinkurl{-tANM6ETl_M$_$3} & 6.758 & 6.604 & 22 & 22 & 21 & 21\\
\nolinkurl{-tPCytz4rww$_$12} & 11.864 & 11.736 & 26 & 26 & 24 & 24\\
\nolinkurl{-tPCytz4rww$_$16} & 7.206 & 7.139 & 16 & 16 & 14 & 14\\
\nolinkurl{-vxjVxOeScU$_$4} & 9.127 & 9.037 & 21 & 21 & 18 & 18\\
\nolinkurl{-wMB_hJL-3o$_$7} & 6.044 & 5.976 & 22 & 22 & 21 & 21\\
\nolinkurl{-wny0OAz3g8$_$3} & 6.146 & 6.100 & 20 & 1 & 20 & 1\\
\nolinkurl{-wny0OAz3g8$_$5} & 6.912 & 6.741 & 20 & 1 & 20 & 1\\
\nolinkurl{-wny0OAz3g8$_$7} & 8.068 & 7.910 & 29 & 6 & 29 & 6\\
\nolinkurl{-wny0OAz3g8$_$9} & 6.213 & 6.043 & 20 & 20 & 19 & 19\\
\addlinespace[0.5em]
\multicolumn{7}{l}{\textbf{3. 全部词通过对齐筛选，但视觉模态不可观测（1条）}}\\*
\nolinkurl{-wny0OAz3g8$_$1} & 6.844 & 6.777 & 22 & 0 & 22 & 0\\
\addlinespace[0.5em]
\multicolumn{7}{l}{\textbf{4. 整段对齐受限（已完成人工审核）（21条）}}\\*
\nolinkurl{-HwX2H8Z4hY$_$9} & 7.734 & 7.636 & 22 & 0 & 0 & 0\\
\nolinkurl{-NFrJFQijFE$_$1} & 5.746 & 5.659 & 16 & 0 & 0 & 0\\
\nolinkurl{-NFrJFQijFE$_$2} & 6.856 & 6.698 & 18 & 0 & 0 & 0\\
\nolinkurl{-hnBHBN8p5A$_$6} & 6.401 & 6.333 & 13 & 0 & 0 & 0\\
\nolinkurl{-hnBHBN8p5A$_$7} & 7.372 & 7.309 & 13 & 0 & 0 & 0\\
\nolinkurl{-iRBcNs9oI8$_$3} & 5.621 & 5.538 & 8 & 0 & 0 & 0\\
\nolinkurl{-iRBcNs9oI8$_$6} & 2.903 & 2.820 & 5 & 0 & 0 & 0\\
\nolinkurl{-iRBcNs9oI8$_$7} & 4.104 & 4.006 & 9 & 0 & 0 & 0\\
\nolinkurl{-iRBcNs9oI8$_$8} & 8.093 & 8.026 & 21 & 0 & 0 & 0\\
\nolinkurl{-iRBcNs9oI8$_$9} & 3.416 & 3.374 & 5 & 0 & 0 & 0\\
\nolinkurl{-mJ2ud6oKI8$_$1} & 6.103 & 6.054 & 13 & 0 & 0 & 0\\
\nolinkurl{-mJ2ud6oKI8$_$2} & 5.731 & 5.702 & 11 & 0 & 0 & 0\\
\nolinkurl{-mJ2ud6oKI8$_$6} & 2.257 & 2.228 & 5 & 5 & 0 & 0\\
\nolinkurl{-mJ2ud6oKI8$_$8} & 4.192 & 4.171 & 11 & 11 & 0 & 0\\
\nolinkurl{-mJ2ud6oKI8$_$9} & 7.033 & 6.909 & 22 & 22 & 0 & 0\\
\nolinkurl{-ri04Z7vwnc$_$0} & 2.807 & 2.763 & 8 & 0 & 0 & 0\\
\nolinkurl{-ri04Z7vwnc$_$2} & 6.027 & 5.913 & 16 & 8 & 0 & 0\\
\nolinkurl{-ri04Z7vwnc$_$5} & 3.545 & 3.440 & 8 & 8 & 0 & 0\\
\nolinkurl{-yRb-Jum7EQ$_$1} & 29.288 & 29.110 & 49 & 48 & 0 & 0\\
\nolinkurl{-yRb-Jum7EQ$_$5} & 13.169 & 13.045 & 25 & 25 & 0 & 0\\
\nolinkurl{-yRb-Jum7EQ$_$6} & 11.243 & 11.243 & 19 & 19 & 0 & 0\\
\end{longtable}
\endgroup

% ===== tab:p1_files =====
\begin{table}[H]
    \centering\small
    \caption{特征结果文件及其内容}
    \label{tab:p1_files}
    \setlength{\tabcolsep}{6pt}
    \renewcommand{\arraystretch}{1.12}
    \ifdefined\pOneCompactTableBox\else\newsavebox{\pOneCompactTableBox}\fi
    \begin{lrbox}{\pOneCompactTableBox}
    \begin{tabular}{l>{\raggedright\arraybackslash}p{24em}}
        \toprule
        文件 & 内容\\
        \midrule
        features.npz & 词级三模态矩阵、候选词区间、观测标记，以及原生视听特征和聚合权重。\\
        metadata.json & 原文、标签、词与子词位置、各维特征名称及处理参数。\\
        \nolinkurl{quality_masks.npz}、\newline quality.json & 对齐质量筛选结果、结合观测状态的有效位置标记及自动筛选原因。\\
        \nolinkurl{windows_0.2s.npz}、\newline\nolinkurl{windows_0.5s.npz} & 固定窗口边界、视听特征及文本聚合权重。\\
        人工审核\_21条样本.csv & 21条样本的人工审核状态、问题分类与原始说明。\\
        三份CSV汇总表 & \nolinkurl{sample_summary.csv}、\nolinkurl{modality_summary.csv}和\nolinkurl{word_alignment.csv}分别对应100条样本、300条样本--模态记录及1932个词的明细。\\
        \bottomrule
    \end{tabular}
    \end{lrbox}
    \usebox{\pOneCompactTableBox}
\end{table}

% ===== tab:p1_repro_settings =====
\begin{table}[H]
    \centering\small
    \caption{实验环境与主要参数}
    \label{tab:p1_repro_settings}
    \setlength{\tabcolsep}{6pt}
    \renewcommand{\arraystretch}{1.12}
    \ifdefined\pOneCompactTableBox\else\newsavebox{\pOneCompactTableBox}\fi
    \begin{lrbox}{\pOneCompactTableBox}
    \begin{tabular}{l>{\raggedright\arraybackslash}p{27em}}
        \toprule
        项目 & 版本或参数\\
        \midrule
        运行环境 & Python 3.10.13；PyTorch 2.5.1；CUDA 12.1。\\
        音视频读取 & PyAV 17.0.1；音频16 kHz、单声道；视频目标抽样频率5 Hz。\\
        文本编码 & Transformers 4.48.3；\nolinkurl{distilbert/distilbert-base-uncased}；每块最多510个子词，另加首尾标记。\\
        声学对齐 & \nolinkurl{facebook/wav2vec2-base-960h}；CTC最大概率路径。\\
        声学特征 & openSMILE 2.6.0；eGeMAPSv02的25维低层描述量。\\
        面部特征 & MediaPipe 0.10.21；Face Landmarker float16版本1；52维表情系数；检测与存在阈值均为0.5，最多两张脸，选择面积最大者。\\
        数值计算 & NumPy 1.26.4；SciPy 1.15.3；特征采用32位浮点数；num2words 0.5.14；随机种子20260923。\\
        \bottomrule
    \end{tabular}
    \end{lrbox}
    \usebox{\pOneCompactTableBox}
\end{table}
